\documentclass[a4paper]{article}
\usepackage[pdftex]{graphicx}
\usepackage[utf8]{inputenc}
\usepackage{enumerate}
\usepackage{siunitx}
\sisetup{locale=DE}
\usepackage{href-ul}
\hypersetup{
	colorlinks=true,
	linkcolor=blue,
	urlcolor=blue}
\usepackage{icomma}
\usepackage{amssymb}
\usepackage{geometry}
\geometry{a4paper, top=15mm, left=15mm, right=15mm, bottom=15mm,
	headsep=10mm, footskip=12mm}

\begin{document}
	{\bf Gravity, Celestial Mechanics}
	\begin{enumerate}[1.]
		
		% Exercise 2
		\begin{minipage}{0.7\textwidth}
			{\bf \item The Martian moon Phobos has a mass of
				$\SI{1,1e16}{\kilogram}$ and a radius of about $\SI{11}{\kilo\meter}$\\}
			$[\textnormal{gravitational constant}~\gamma=\SI{6.67e-11}{\meter^3 \per {\kilogram\second^2}}]$
		\end{minipage}
		\hfill
		\begin{minipage}{0.25\textwidth}
			\includegraphics[width=2.5 cm]{phobos95}
		\end{minipage}
		
		\begin{enumerate}[a)]
			\item Show by calculation that there is a gravitational acceleration of approximately $g_{Ph}=\SI{6,1e-3}{\meter\over {\second^2}}$ on the surface of Phobos.
			\item On Phobos you can play baseball with yourself. To do this, you throw a ball horizontally so that it circles Phobos in a circular path close to the surface and arrives back at the throwing point. At what speed does this ball have to be thrown?
			\item After the ball is thrown, wait until the ball reappears on the opposite horizon. How much time passes from the throw to the completion of the circuit, i.e. to the tee shot?
		\end{enumerate}
		
		
		{\bf \item On Earth the location factor $g$ and the radius} depend on the geographical latitude. The radius of the Earth is $\SI{6357}{\kilo\meter}$ at the pole, $\SI{6378}{\kilo\meter}$ at the equator, the mass of the Earth is $\SI{5.97e24}{\ kilogram}$
		
		\begin{enumerate}[a)]
			\item Which two effects are responsible for the difference between $g_{Pol}-g_{Aeq}$?
			\item Calculate the absolute and relative difference between $g_{Pol}$ and $g_{Aeq}$
		\end{enumerate}
		
		{\bf \item Black holes have such a strong gravitational field} that not even light is fast enough to escape if it is closer than the Schwarzschild radius. Calculate this radius for a black hole with 10 times the mass of the Sun $M_{Sun}= \SI{1.98e30}{\kilogram}$
		
		{\bf \item There is a point between the moon and the earth} in which the gravitational field strengths cancel each other out. Calculate the distance $r$ of this point to the center of the moon. Let the Moon-Earth distance be $d=\SI{384400}{\kilo\meter}$. Mass of the Moon: $\SI{7.35e22}{\kilogram}$
		
		\includegraphics[width=18 cm]{moe413}
		{\bf \item Spaceman Spiff swung his lander into orbit around an exoplanet}.
		\noindent
		\begin{minipage}{0.4\textwidth}
			It measures the following data:
			\begin{enumerate}[a)]
				\item What is the radius of the planet?
				\item What is the mass of the planet?
				\item What gravitational field strength will Spiff expect when landing?
			\end{enumerate}
			
		\end{minipage}
		\hspace{1cm}
		\begin{minipage}{0.4\textwidth}
			\renewcommand{\arraystretch}{1.5}
			\begin{tabular}{|c|c|}
				\hline
				Height above ground & $h=\SI{800}{\kilo\meter}$ \\
				\hline
				Speed over ground & $v=\SI{3,7}{\kilo\meter \over \second}$\\
				\hline
				Orbital period & $T=\SI{130}{\minute}$\\
				\hline
			\end{tabular}
		\end{minipage}
		
		[The exoplanet's own rotation is negligible!]
		
		
		
		
		
	\end{enumerate}
	
	\fbox{
		\begin{minipage}{0.5\textwidth}
			For the solution, please \href{https://www.okuyakl.de/phys/p10kreiL413/le413.pdf}{click here} or scan the QR code.\\
			Additional worksheets are available at
			
			\href{http://www.okuyakl.com}{www.okuyakl.com}
		\end{minipage}
		\hfill
		\begin{minipage}{0.4\textwidth}
			\includegraphics[width=1.5 cm]{../../viecher/zwe03}
			\includegraphics[width=3 cm]{qre413}
			\includegraphics[width=2 cm]{../../viecher/afanticon1}
			
	\end{minipage}}
	
\end{document}